\subsection{Some more recent achievements}
\label{sec:recent-achievements}
\slidesrc{Thinking machines: 68--79}

After 2016 the pace accelerated. The slides show a selection of results; for each one,
here is what it is and why it matters.

\subsubsection{Games, medicine and emergent behaviour}
\slidesrc{Thinking machines: 68--71}

\begin{itemize}
  \item \textbf{AlphaStar} (DeepMind, 2019). StarCraft~II is a real-time strategy video game:
        many units, partial information (you do not see the whole map), real-time decisions.
        AlphaStar won its historic series against professional players (announced on
        January 24, 2019, after matches recorded in December 2018; e.g.\ 5--0 against the pro
        player MaNa).
  \item \textbf{Lung cancer screening} (\emph{Nature Medicine}, May 2019). A
        three-dimensional deep learning model analyses low-dose chest CT scans end to end and
        detects lung cancer with performance comparable to, or better than, expert
        radiologists.
  \item \textbf{Emergent tool use from multi-agent autocurricula} (OpenAI, ICLR 2020). Agents
        play \emph{hide-and-seek} in a simulated world with movable boxes and ramps. Nobody
        tells them how to use the objects; just by competing, the teams discover a sequence
        of strategies: hiders build shelters with boxes, seekers use ramps to jump in,
        hiders then lock the ramps away, seekers ``surf'' on boxes, and so on.
\end{itemize}

\begin{nbwhy}[A multi-agent system in action]
The hide-and-seek experiment is a \emph{multi-agent system}: several autonomous agents,
in a shared environment, competing (hiders vs.\ seekers) and cooperating (inside a team).
Complex behaviour \emph{emerges} from simple goals and interaction, with no central
controller. Keep this example in mind for Part~III (agents and MAS) and for agent-based
simulation (Section~\ref{sec:abms}).
\end{nbwhy}

\subsubsection{Self-driving cars: the SAE levels of driving automation}
\slidesrc{Thinking machines: 72--74}

How ``autonomous'' is an autonomous car? The SAE J3016 standard defines six levels.

\begin{center}
\small
\renewcommand{\arraystretch}{1.15}
\begin{tabularx}{\linewidth}{@{}c >{\raggedright\arraybackslash}X >{\raggedright\arraybackslash}X >{\raggedright\arraybackslash}X@{}}
\toprule
\textbf{Level} & \textbf{What the human must do} & \textbf{What the features do} & \textbf{Examples} \\
\midrule
\rowcolor{nbblue!6} 0 & You \textbf{are driving}, even if features are on; constantly supervise & warnings and momentary assistance & automatic emergency braking, blind-spot warning, lane departure warning \\
\rowcolor{nbblue!6} 1 & same as level 0 & steering \textbf{or} brake/acceleration support & lane centering \textbf{or} adaptive cruise control \\
\rowcolor{nbblue!6} 2 & same as level 0 & steering \textbf{and} brake/acceleration support & lane centering \textbf{and} adaptive cruise control together \\
\rowcolor{nbgreen!8} 3 & You are \textbf{not driving} while features are engaged, but you must take over when the system requests it & drives the vehicle under limited conditions & traffic jam chauffeur \\
\rowcolor{nbgreen!8} 4 & You are not driving; you will never be required to take over & drives under limited conditions, never needs you & local driverless taxi; pedals and steering wheel may not be installed \\
\rowcolor{nbgreen!8} 5 & You are not driving & drives everywhere, in all conditions & as level 4, but anywhere \\
\bottomrule
\end{tabularx}
\captionof{table}{SAE levels of driving automation (slide 72). Blue rows: \emph{driver support}
features (the human drives). Green rows: \emph{automated driving} features (the system
drives).}
\end{center}

The slides note (2024 news) that a \textbf{Level~3} car can now be bought in the US: the
Mercedes-Benz system works only on a small selection of roads, at limited speed and only in
some states, and its wide adoption will require regulators to catch up.

\begin{nbintuition}[Levels of autonomy, again later]
The question ``how much can the system decide by itself, and when must a human step in?''
is the question of \textbf{autonomy}, the defining property of agents
(Section~\ref{sec:agents-and-mas}). The jump from level 2 to level 3 is exactly the moment
in which responsibility moves from the human to the agent.
\end{nbintuition}

\subsubsection{Generative AI, protein folding and agentic AI}
\slidesrc{Thinking machines: 75--79}

On slide 75 the professor asked ChatGPT itself: ``which are the three most impressive
achievements in the AI research field in the last three years (short answer)''. The answer:
\begin{enumerate}[nosep]
  \item \textbf{Generative AI}: models such as GPT-4 and DALL-E generate coherent text,
        images and creative content.
  \item \textbf{AI in protein folding}: DeepMind's AlphaFold accurately predicts protein
        structures, solving a 50-year-old problem in biology.
  \item \textbf{AI alignment and safety research}: progress in interpretability, alignment
        and safety (e.g.\ by OpenAI and Anthropic) to control the behaviour of advanced AI.
\end{enumerate}

\begin{nbdef}{Protein folding problem}
A protein is a chain of amino acids. The \textbf{protein folding problem} is the question of
how a protein's amino acid sequence dictates its three-dimensional atomic structure. It
first emerged around 1960, with the first atomic-resolution protein structures.
\end{nbdef}

\textbf{AlphaFold2} (DeepMind) predicts the 3D structure of proteins from their amino acid
sequence with near-experimental accuracy; more than 200 million predicted structures have
been released. Knowing a protein's shape tells biologists how it works, which speeds up
drug discovery and protein design (2023 review; PNAS 2024, ``AlphaFold two years on'').

Finally, slide 79: \textbf{agentic AI}, presented as ``the next evolution of artificial
intelligence''. An agentic AI system does not just answer one question: given a goal, it
plans, calls tools (search engines, APIs, other programs), observes the results and acts
again, autonomously, over many steps.

\begin{nbwhy}[Old ideas, new engines]
Agentic AI is the 2020s name for an idea this course has studied for decades:
\textbf{autonomous agents} that perceive, decide and act towards goals, and
\textbf{multi-agent systems} in which several of them coordinate. The difference is the
engine: a large language model instead of (or together with) explicit beliefs, goals and
plans. Understanding classic agents (BDI, communication protocols, platforms) is what makes
it possible to engineer agentic AI systems that are reliable.
\end{nbwhy}
